
\documentclass{article}
\usepackage{colortbl}
\usepackage{makecell}
\usepackage{multirow}
\usepackage{supertabular}

\begin{document}

\newcounter{utterance}

\centering \large Interaction Transcript for game `function\_detective', experiment `list\_passive\_labeled\_pairs\_oneshot', episode 113 with Qwen3.6{-}35B{-}A3B{-}FP8{-}without{-}reasoning.
\vspace{24pt}

{ \footnotesize  \setcounter{utterance}{1}
\setlength{\tabcolsep}{0pt}
\begin{supertabular}{c@{$\;$}|p{.15\linewidth}@{}p{.15\linewidth}p{.15\linewidth}p{.15\linewidth}p{.15\linewidth}p{.15\linewidth}}
    \# & \multicolumn{2}{c}{Player} && \multicolumn{2}{c}{Game Master} \\
    \hline

    \theutterance \stepcounter{utterance}  
    & & & \multicolumn{4}{p{0.6\linewidth}}{
        \cellcolor[rgb]{0.9,0.9,0.9}{
            \makecell[{{p{\linewidth}}}]{
                \texttt{\tiny{[P1$\langle$GM]}}
                \texttt{I have a hidden Python function for you to figure out.} \\
\\ 
\texttt{Your job is to infer what it does and then submit a correct Python implementation.} \\
\\ 
\texttt{Function signature:} \\
\texttt{(items: List{[}int{]}) {-}> int} \\
\\ 
\texttt{You have 1 turns total.} \\
\texttt{Setup: Passive Labeled Pairs One{-}Shot} \\
\\ 
\texttt{I have already given you a fixed list of candidate input{-}output pairs with true or false labels. Solve using those statements only.} \\
\\ 
\texttt{You must make your best SOLVE attempt no later than the second{-}to{-}last turn.} \\
\texttt{Do not wait until the very end.} \\
\\ 
\texttt{You can use these actions:} \\
\texttt{{-} SOLVE: ```python ... ```} \\
\\ 
\texttt{Final answer format:} \\
\texttt{SOLVE: ```python} \\
\texttt{def solution(<value1>):} \\
\texttt{    ...} \\
\texttt{```} \\
\\ 
\texttt{Do not add import statements in your guess.} \\
\texttt{Anything you might need is already available, so just write the function body and return the answer.} \\
\\ 
\texttt{No further example requests are allowed in this one{-}shot baseline.} \\
\\ 
\texttt{I will reply with OUTPUT: lines only.} \\
\texttt{You must never write OUTPUT: yourself.} \\
\\ 
\texttt{Preloaded examples:} \\
\texttt{1. (({[}{-}1, 0, 1{]}, 1), True)} \\
\texttt{2. (({[}1, 2, 3{]}, {-}1), False)} \\
\texttt{3. (({[}1, 1{]}, 0), True)} \\
\texttt{4. (({[}0{]}, 2), False)} \\
\texttt{5. (({[}2, 1{]}, 0), True)} \\
\texttt{6. (({[}1{]}, {-}1), False)} \\
\texttt{7. (({[}5, {-}5, 5{]}, 0), True)} \\
\texttt{8. (({[}{]}, {-}1), False)} \\
\texttt{9. (({[}3, 2, 1{]}, 0), True)} \\
\texttt{10. (({[}1, {-}10, 10, 0, {-}2{]}, 2), False)} \\
\\ 
\texttt{Scoring:} \\
\texttt{{-} correctness is required for any non{-}zero score} \\
\texttt{{-} in interactive modes, solving with fewer information requests is better} \\
\\ 
\texttt{Input formatting reminders:} \\
\texttt{{-} strings must be quoted} \\
\texttt{{-} booleans must be True/False} \\
\texttt{{-} use commas between multiple values} \\
\texttt{{-} do not write any extra action tags in the same message} \\
\\ 
\texttt{Format example:} \\
\texttt{You write:} \\
\texttt{SOLVE: ```python} \\
\texttt{I answer:} \\
\texttt{def solution(...):} \\
\texttt{    ...} \\
\texttt{```} \\
\\ 
\texttt{Output only one valid action per turn.} \\
            }
        }
    }
     \\ \\

    \theutterance \stepcounter{utterance}  
    & \multicolumn{4}{p{0.6\linewidth}}{
        \cellcolor[rgb]{0.561,0.737,0.561}{
            \makecell[{{p{\linewidth}}}]{
                \texttt{\tiny{[P1$\rangle$GM]}}
                \texttt{Let's analyze the input{-}output pairs to infer the function's behavior.} \\
\\ 
\texttt{The function signature is `def solution(items: List{[}int{]}) {-}> int`, but the examples show pairs of `(input\_list, target\_value)` resulting in `True` or `False`. This suggests the hidden function might actually be checking if a certain condition is met, or perhaps the "hidden function" we need to implement is a predicate that takes a list and a target, returning a boolean. However, the prompt asks for a function with signature `(items: List{[}int{]}) {-}> int`. Let's re{-}read carefully.} \\
\\ 
\texttt{"Function signature: (items: List{[}int{]}) {-}> int"} \\
\texttt{But the examples are `((list, int), bool)`.} \\
\\ 
\texttt{This is a contradiction. Let's look at the examples again.} \\
\texttt{1. `({[}{-}1, 0, 1{]}, 1) {-}> True`} \\
\texttt{2. `({[}1, 2, 3{]}, {-}1) {-}> False`} \\
\texttt{3. `({[}1, 1{]}, 0) {-}> True`} \\
\texttt{4. `({[}0{]}, 2) {-}> False`} \\
\texttt{5. `({[}2, 1{]}, 0) {-}> True`} \\
\texttt{6. `({[}1{]}, {-}1) {-}> False`} \\
\texttt{7. `({[}5, {-}5, 5{]}, 0) {-}> True`} \\
\texttt{8. `({[}{]}, {-}1) {-}> False`} \\
\texttt{9. `({[}3, 2, 1{]},} \\
            }
        }
    }
    & & \\ \\

    \theutterance \stepcounter{utterance}  
    & & & \multicolumn{2}{p{0.3\linewidth}}{
        \cellcolor[rgb]{0.9,0.9,0.9}{
            \makecell[{{p{\linewidth}}}]{
                \texttt{\tiny{[GM$|$GM]}}
                \texttt{Game Verdict: LOSS (turn limit reached)} \\
            }
        }
    }
    & & \\ \\

\end{supertabular}
}

\end{document}
